% =============================================================================
\section{Problématique}
% =============================================================================

\begin{frame}{Problématique : objectifs, contraintes et difficultés}
  \footnotesize
  \begin{columns}[T,onlytextwidth]
    \begin{column}{0.48\textwidth}
      \begin{block}{Objectif}
        Affecter chaque patient à une \textbf{vacation} compatible avec sa
        spécialité et à un \textbf{lit}, pour :
        \begin{itemize}
          \item occuper le bloc sans dépasser les vacations ;
          \item lisser les lits et les places ambulatoires ;
          \item réserver une marge pour les urgences.
        \end{itemize}
      \end{block}
      \begin{block}{Contraintes}
        \begin{itemize}
          \item \textbf{Médicales} : spécialité, urgence ;
          \item \textbf{Organisationnelles} : chirurgien, équipes ;
          \item \textbf{Temporelles} : jour, horaires, séjour.
        \end{itemize}
      \end{block}
    \end{column}
    \begin{column}{0.48\textwidth}
      \begin{alertblock}{Difficultés}
        \begin{itemize}
          \item \textbf{Combinatoire élevée} : le produit des affectations
            explose.
          \item \textbf{Multi-critères} : bloc, lits et ambulatoire
            s'opposent.
          \item \textbf{Dépendances fortes} : le bloc retentit sur les lits
            plusieurs jours après.
        \end{itemize}
      \end{alertblock}
      \begin{center}
        Problème d'affectation sous contraintes : \textbf{NP-difficile}.
      \end{center}
    \end{column}
  \end{columns}
  \source{Formulation : affectation patient $\rightarrow$ salle $\rightarrow$
    jour (variables $x_{i,s,t}$).}
\end{frame}

